\documentclass[12pt]{article}

\usepackage[a4paper, margin=1in]{geometry}
\usepackage{array}
\usepackage{longtable}
\usepackage{titlesec}
\usepackage{xcolor}
\usepackage{graphicx}
\usepackage{float}

% Define custom colours
\definecolor{headerblue}{RGB}{30,60,120}
\definecolor{veryminor}{RGB}{0,130,0}
\definecolor{minor}{RGB}{90,140,0}
\definecolor{significant}{RGB}{180,140,0}
\definecolor{major}{RGB}{200,100,0}
\definecolor{severe}{RGB}{180,0,0}
\definecolor{rare}{RGB}{0,130,0}
\definecolor{unlikely}{RGB}{90,140,0}
\definecolor{moderate}{RGB}{180,140,0}
\definecolor{high}{RGB}{200,100,0}
\definecolor{certain}{RGB}{180,0,0}

% Section styling
\titleformat{\section}
  {\large\bfseries\color{headerblue}}
  {}{0em}{}
\setlength{\parindent}{0pt}

\title{ENTER YOUR VULNERABILITY FINDING TITLE HERE}
\author{YOUR NAME}
\date{\today}

\begin{document}

\maketitle

\section*{Assessment Details}

\begin{tabular}{|p{5cm}|p{7.5cm}|}
\hline
Name & ENTER YOUR NAME \\ \hline
Team & App Attack PenTesting \\ \hline
Role & JUNIOR/SENIOR MEMBER/LEAD \\ \hline
Project & PROJECT YOU PEN TESTED \\ \hline
Quality Assurance & Leave blank until QA \\ \hline
Was this Finding Successful? & Yes \\ \hline
Is this a Re-tested Finding? & No \\ \hline
\end{tabular}

\section*{Impact Values}

\begin{longtable}{|p{3cm}|p{12cm}|}
\hline
\textbf{\textcolor{veryminor}{Very Minor}} & Risk that holds little to no impact. Will not cause damage and regular activity can continue. \\ \hline
\textbf{\textcolor{minor}{Minor}} & \textcolor{minor}{Risk that holds minor impact. Can cause some damage but not enough to significantly affect operations.} \\ \hline
\textbf{\textcolor{significant}{Significant}} & Risk that causes noticeable disruption and may impede regular activity. \\ \hline
\textbf{\textcolor{major}{Major}} & Risk that significantly impacts operations and prevents the system from running normally. \\ \hline
\textbf{\textcolor{severe}{Severe}} & Risk that causes critical damage and may stop activity entirely. \\ \hline
\end{longtable}

\section*{Likelihood}

\begin{longtable}{|p{3cm}|p{12cm}|}
\hline
\textbf{\textcolor{rare}{Rare}} & Event may occur only in exceptional circumstances. \\ \hline
\textbf{\textcolor{unlikely}{Unlikely}} & Event could occur occasionally. \\ \hline
\textbf{\textcolor{moderate}{Moderate}} & \textcolor{moderate}{Event may occur under certain conditions.} \\ \hline
\textbf{\textcolor{high}{High}} & Event is likely to occur often. \\ \hline
\textbf{\textcolor{certain}{Certain}} & Event is occurring now or is expected to occur frequently. \\ \hline
\end{longtable}

\section*{Finding Details}

\subsection*{Finding Name: ENTER HERE}

\subsubsection*{Description}
EXPLAIN VULNERABILITY, AND WHY IT MATTERS

\subsubsection*{Risk}
EXPLAIN WHAT AN ATTACKER COULD ACHIEVE.

\subsubsection*{Affected Assets}
\begin{itemize}
    \item AFFECTED ASSETS
\end{itemize}

\subsubsection*{Affected Path}
AFFECTED URL/FILEPATH

\subsubsection*{Evidence}

\vspace{0.5em}

INCLUDE DESCRIPTIONS ALONGSIDE IMAGE EVIDENCE

\begin{figure}[H]
\includegraphics[width=1\linewidth]{IMAGE NAME OR FILEPATH RELATIVE TO LATEX DOC.}
\caption{IMAGE CAPTION}
\end{figure}

\vspace{0.5em}

\subsubsection*{Remediation Advice}
\begin{itemize}
    \item REMEDIATION RECOMMENDATIONS
\end{itemize}
\end{document}
